\documentclass{beamer}
%  Use something like this Perl command to substitute across all files:
% $ perl -p -i -e 's/\\subfolder\//\\dayonesegone\//g' *.tex

% Use jpg graphics in figure environments and pdflatex to compile!!!
\usepackage{graphicx,color,amsmath,amssymb,amsfonts,amsbsy,amsthm}
\usepackage{verbatim}
\usepackage{subfig}
\usetheme{Warsaw}
\newcommand{\includes}{/cygwin64/home/cognotrend/kalman_course/kf_gps_v4/include}
\newcommand{\graphics}{/cygwin64/home/cognotrend/kalman_course/kf_gps_v4/graphics}
%\newcommand{\includes}{/home/1104070/latex_include}
\newcommand{\dayonesegone}{.}
% Node the following tag will orient things right using dvi2pdf.
% Process with LaTeX first, then dvi2pdf.
\special{landscape}
% Fuzz -----------------------------------------------------------
\hfuzz5pt % Don't bother to report overfull boxes < 5pt
% MATH -----------------------------------------------------------
\input{\includes/kalman_cmdmaster.tex}
\input{\includes/orbitfilter_cmds.tex}
% THEOREM ENV. ---------------------------------------------------
\newtheorem{thm}{Theorem.}
\newtheorem{cor}[thm]{Corollary.}
\newtheorem{con}[thm]{Conjecture.}
\newtheorem{lem}[thm]{Lemma.}
\newtheorem{prop}[thm]{Proposition.}
\newtheorem{defn}[thm]{Definition.}
\renewcommand{\thethm}{}  % unnumbered
\pagestyle{plain}
%%% --------------------------------------------------------------
\newcounter{exercisecounter}

\begin{document}
\title{
\vfill \hrulefill \vfill
\begin{center}
Orbital Perturbations: Variation of Parameters
\end{center}
}
\author{Michael L. Carroll}
\date{\today\\
\vfill \hrulefill \vspace*{.25in} \vfill
\begin{tiny}
\copyright 2020 by Michael L. Carroll
\end{tiny}
}
\maketitle

% ----------------------------------------------------------------
%\input{equations_main_video4.tex}
\begin{frame}
  \onslide<1->{
	\frametitle{Catalog of Orbital Equations} 
	\framesubtitle{}
}
\begin{itemize}
  \onslide<2->{
   \item Basic Two-Body Equation: $\twobodyeq$,\\ 
      $\twobodywhere$ 
}
  \onslide<3->{
   \item $\mu$ of Earth is $\muearth$
}
  \onslide<4->{
   \item Angular Momentum:  $\angmom$
}
  \onslide<5->{
   \item Classical Orbital Elements: $a,e,i,\Omega,\omega,T$ where\\
   $\semimajordef$,\\
   $\eccdef$,\\
   $\incdef$,\\
   $\longanode$,\\
   $\argperiapsis$,\\
   $\timeperiapsis$.
}
  \onslide<6->{
   \item
}
  \onslide<7->{
   \item
}
  \onslide<8->{
   \item
}
  \onslide<9->{
  \item

}
\end{itemize}
\onslide<10->{\center --------------------------------------------------------------------}
\end{frame}

\end{document}
